% ===========================================================================
\section{Conclusiones}
\label{sec:conclusiones}

Se implementó y verificó en la Basys~3 un procesador RISC-V RV32I segmentado en cinco etapas, con
forwarding, detección de carga-uso y flush por saltos, controlado desde la PC por una Debug Unit que
carga programas, ejecuta de corrido o de a un ciclo y vuelca registros, latches y memoria de datos
usada. Funciona a \qty{92}{\mega\hertz}, la máxima frecuencia que cumple timing, con un solo reloj y
sin intervenirlo.

\paragraph{Lo que funcionó.}
\begin{itemize}[nosep]
  \item \textbf{Especificar antes de codificar.} Las interfaces de cada módulo y el contrato de
        nombres de los registros de segmentación se fijaron en \repo{docs/spec/} antes del RTL. El
        mismo empaquetado de los latches sirve para la traza del modelo, el banco de pruebas y el
        volcado de la Debug Unit, así que la placa se compara con el modelo sin traducciones.
  \item \textbf{Integrar etapa por etapa.} Cada error quedó acotado a la última etapa agregada.
  \item \textbf{Dos oráculos independientes y mutantes.} Con los \texttt{.exp} a mano y el
        modelo escritos por separado, cada diferencia señalaba qué revisar (la ISA, el modelo o el
        RTL); los mutantes mostraron qué programas faltaban para cubrir cada mecanismo.
  \item \textbf{Medir antes de optimizar.} La hipótesis del camino crítico de la \decision{012} se
        confirmó en el qué y se corrigió en el cómo (el desplazador, no la resta), y eso llevó a una
        mejora que no cambió ciclos ni resultados: el comparador de igualdad.
\end{itemize}

\paragraph{Limitaciones y trabajo futuro.}
\begin{itemize}[nosep]
  \item Sin excepciones: los accesos inválidos se truncan y dejan un aviso (\decision{015}).
  \item Los saltos tomados cuestan 2 ciclos. Una predicción de saltos hacia atrás o el forwarding
        hacia MEM para un \texttt{sw} después de un \texttt{lw} ahorrarían ciclos a cambio de más
        lógica en el camino crítico.
  \item A \qty{92}{\mega\hertz} varias etapas quedan en el límite y todas empiezan en la BRAM de
        datos; subir más exige registrar su salida o la redirección, con un ciclo más por load o por
        salto.
\end{itemize}
